% =========================================================
% Proof Stub: Absolute Convergence Implies Convergence
% Source: volume-iii/analysis/sequences/notes/series/notes-absolute-convergence.tex
% Mode: standard
% =========================================================

\clearpage
\phantomsection
\hypertarget{proof-RA-ABB-T-abs-conv-implies-conv}{}

\subsubsection[Absolute Convergence Implies Convergence]{Proof --- Absolute Convergence Implies Convergence}

\bigskip

\noindent
\textbf{Source.}
Abbott, \emph{Understanding Analysis}; see \texttt{notes-absolute-convergence.tex}.

\vspace{0.75em}

\noindent
\textbf{Goal.}
If $\sum |a_n| < \infty$, then $\sum a_n$ converges.

\vspace{0.75em}

\noindent
\textbf{Logical form.}
\[
  \sum |a_n| < \infty \Rightarrow \sum a_n \text{ converges}.
\]

\vspace{0.75em}

\noindent
\textbf{Proof strategy.}
Use the Cauchy Criterion: $|s_m - s_n| = |\sum_{k=n+1}^m a_k| \le \sum_{k=n+1}^m |a_k|$. Since $\sum |a_n|$ converges, its tail sums go to 0, making $(s_n)$ Cauchy.

\vspace{0.75em}

\noindent
\textbf{Proof.}
\begin{proof}
\Claim If $\sum |a_n| < \infty$, then $\sum a_n$ converges.

\Given 

\Goal 

\Strategy Use the Cauchy Criterion: $|s_m - s_n| = |\sum_{k=n+1}^m a_k| \le \sum_{k=n+1}^m |a_k|$. Since $\sum |a_n|$ converges, its tail sums go to 0, making $(s_n)$ Cauchy.

\medskip

% --- |s_m - s_n| = |∑_{k=n+1}^m a_k| ≤ ∑_{k=n+1}^m |a_k| ---
% --- ∑|a_n| converges ⟹ its partial sums are Cauchy ---
% --- so ∑_{k=n+1}^m |a_k| < ε for m > n ≥ N ---
% --- therefore (s_n) is Cauchy; Cauchy Criterion ⟹ (s_n) converges ---

\end{proof}

\begin{remark*}[Proof shape]
Cauchy Criterion: bound $|s_m - s_n|$ by the tail of $\sum |a_n|$, which goes to 0.
\end{remark*}

\begin{remark*}[Dependencies]
\begin{itemize}
  \item Cauchy Criterion in $\mathbb{R}$.
  \item Triangle inequality for finite sums.
  \item Convergent series has Cauchy partial sums.
\end{itemize}
\end{remark*}
